% CV 2026 -- single-column, ATS-friendly template (xelatex).
\documentclass[a4paper,11pt]{article}

% ---------------------------------------------------------------- Settings
\newcommand{\cvname}{Adrien Lemaire}
\newcommand{\cvheadline}{CTO \& Co-founder · Full-stack Engineer}

% ---------------------------------------------------------------- Packages
\usepackage{fontspec}
\usepackage[margin=2.2cm]{geometry}
\usepackage{xcolor}
\usepackage{enumitem}
\usepackage{needspace}
\usepackage{etoolbox}
\usepackage{hyperref}

\setmainfont{Linux Libertine O}[Ligatures={TeX,NoCommon}]
\setsansfont{Linux Biolinum O}[Ligatures={TeX,NoCommon}]

\definecolor{accent}{RGB}{31,78,121}
\definecolor{muted}{RGB}{90,90,90}

\hypersetup{
  pdftitle={\cvname{} -- CV},
  pdfauthor={\cvname},
  pdfsubject={Curriculum Vitae},
  colorlinks=true,
  allcolors=accent,
}

\pagestyle{empty}
\linespread{1.02}
\setlength{\parindent}{0pt}
\setlength{\parskip}{0pt}
\setlength{\emergencystretch}{2em}
\hyphenpenalty=10000 % keep keywords whole for ATS parsers
\setlist[itemize]{leftmargin=1.2em, label=\textbullet, topsep=4pt,
  itemsep=3pt, parsep=0pt}

% ---------------------------------------------------------------- Macros
% \cvsection{title}
\newcommand{\cvsection}[1]{%
  \par\bigskip\needspace{4\baselineskip}%
  {\sffamily\large\bfseries\color{accent}#1}\par
  \smallskip{\color{accent}\hrule height 0.6pt}\addvspace{\medskipamount}}

% \role{company}{title}{dates}{location}
\newcommand{\role}[4]{%
  \par\needspace{5\baselineskip}\addvspace{\bigskipamount}
  {\bfseries #1} -- #2\ifstrempty{#3}{}{\hfill #3}\par
  {\small\color{muted}#4}\par\smallskip}

% \roledesc{one-line description of the company}
\newcommand{\roledesc}[1]{{\small\itshape #1}\par}

% \stack{comma-separated technologies}
\newcommand{\stack}[1]{{\small\textbf{Stack:} #1}\par}

% \entry{label}{text}
\newcommand{\entry}[2]{\textbf{#1:} #2}

% ================================================================= Document
\begin{document}

{\sffamily\Huge\bfseries\color{accent}\cvname}\par
\smallskip
{\large\cvheadline}\par
\smallskip
{\raggedright Toulouse, France
· \href{mailto:lemaire.adrien@gmail.com}{lemaire.adrien@gmail.com}
· \href{https://adrienlemaire.github.io/cv}{adrienlemaire.github.io/cv}
· \href{https://github.com/AdrienLemaire}{github.com/AdrienLemaire}
· \href{https://linkedin.com/in/adrienlemaire}{linkedin.com/in/adrienlemaire}\par}

\cvsection{Summary}
CTO and full-stack engineer with 16 years of experience taking products from
idea to production in France, Japan and the UK. Co-founded Made In Tracker, a
deeptech startup, and built MD Atlas, a medical-device regulatory search
engine with 110,000+ users across 170 countries. Hands-on across architecture,
data pipelines, AI tooling and team leadership.

\cvsection{Experience}

\role{Made In Tracker}{Co-founder \& CTO}{2023/03 -- present}{Toulouse, France (remote)}
\roledesc{Traceability and regulatory data software for the medical-device
  industry. Labelled deeptech startup; €90k from angel investors and €90k in
  public funding; 1 patent; valued at €1.7M in 2026.}
\begin{itemize}
  \item Built MD Atlas (\href{https://md-atlas.com}{md-atlas.com}), a
    regulatory intelligence search engine for medical devices: 5.2M records
    across companies, devices, certificates, people and documents, with
    conformity alerts and feed monitoring; 110,000+ users across 170 countries.
  \item Shipped its AI layer, including an MCP server listed in the official
    MCP Registry, so AI assistants can query European regulatory data
    directly.
  \item Led EasyUDI, a SaaS for UDI and EUDAMED compliance: GS1 label and
    DataMatrix generation, XML validation, traceability, audit trail and an
    AI chatbot.
  \item Took R\&D to pilots: product passports, smart labels and secure QR
    codes; Artwork Passport (AI-assisted registration); pilots with Anopa
    (cotton) and SobraBTP (construction reuse). 1st place at the Axone
    blockchain contest.
  \item Grew the team to 8 people at once (about 20 over four years) and
    coached 12+ developers. Owned architecture, stack and conventions, code
    review, infrastructure and client integrations; worked with design, sales
    and investors.
\end{itemize}
\stack{TypeScript, React, Flutter, Python, PostgreSQL/Supabase, Firebase/GCP,
  Vercel, Docker, Nix, LLM/RAG, MCP}

\role{YAOYA Inc.}{CTO}{2017/05 -- 2023/02}{Osaka, Japan / remote}
\begin{itemize}
  \item Built three versions of GOunite (internship marketplace, then online
    school, then metaverse community) on GRANDstack (React, Apollo,
    TypeScript, Neo4j) and Pixi.js, reaching 2,500+ users; started a Unity
    version in 2022.
  \item Rewrote Reach@, the company's article-writing service, from
    PHP/jQuery to Django Wagtail, MySQL and Kubernetes: 2,000+ articles a
    month, ¥100M+ yearly revenue.
  \item Led a team of about 4 developers plus interns, and supported other
    divisions with automation, data backups and client website audits.
\end{itemize}

\role{ShareWis Inc.}{Senior Full-stack Developer}{2016/02 -- 2016/12}{Osaka, Japan}
\begin{itemize}
  \item Technical lead on EdTech products: the ShareWis ACT MOOC marketplace
    (Ruby on Rails, i18n, payments), React Native apps for ShareWis ACT and
    ShareWis snack (authentication, payments, premium features), and a
    Serverless account API.
  \item Built a Node.js automation linking AWS, Pivotal Tracker, GitHub and
    Slack; led the internationalisation effort.
\end{itemize}

\role{Rakuten}{Software Engineer}{2015/09 -- 2016/02}{Osaka, Japan}
\begin{itemize}
  \item Started the large-scale PHP/Symfony migration of O-net, Japan's
    largest marriage agency (43,500+ members), and ran a security audit.
\end{itemize}

\role{Alpaca.ai (formerly Ikkyo Technology)}{Engineer}{2014/06 -- 2015/09}{Kobe, Japan}
\begin{itemize}
  \item Core design of the AlpacaDB platform, released as Labellio, a
    deep-learning image labelling and training service (Kubernetes, Docker,
    Django, Caffe, RabbitMQ).
  \item Computer-vision R\&D for major Japanese gaming and search companies;
    backend of Iromo, a fashion search engine matching clothes by colour and
    pattern.
\end{itemize}

\role{Aquasys}{Director of Technology (2012/11 -- 2014/08),
  Software Engineer (2011/04 -- 2012/11)}{}{Okinawa, Japan}
\begin{itemize}
  \item Supervised all technical operations; hired, trained and managed
    freelancers; built a Django marketplace framework (Uruma-work) and ran
    Odoo operations.
\end{itemize}

\role{Fry-IT}{Django Developer}{2010/07 -- 2011/03}{London, UK}
\begin{itemize}
  \item Built a Django e-portfolio for London hospitals (PostgreSQL, Celery,
    RabbitMQ) using TDD.
\end{itemize}

\cvsection{Skills}
\begin{itemize}
  \item \entry{Leadership}{technical strategy, architecture, hiring and
    coaching, product scoping, investor relations}
  \item \entry{Languages \& frameworks}{TypeScript, JavaScript, Python, Dart;
    React, Node.js, Flutter, Django}
  \item \entry{Data \& AI}{PostgreSQL, Supabase, Neo4j, data pipelines,
    LLM/RAG, MCP, computer vision}
  \item \entry{Infrastructure}{GCP, Firebase, Vercel, AWS, Docker, Kubernetes,
    Nix, CI/CD}
  \item \entry{Practices}{TDD, DDD, Lean Startup, code review}
\end{itemize}

\cvsection{Languages}
French (native) · English (bilingual) · Japanese (professional)

\cvsection{Education}
\textbf{SUPINFO International University} -- International Master of Science
in Computer Science\hfill 2006 -- 2011\par
{\small\color{muted}Studied in France, Canada, the UK and the USA}

\cvsection{Certifications}
Wharton Entrepreneurship specialisation (2016) · Stanford Machine Learning (2011)

\cvsection{Interests}
Triathlon, ultra-marathon

\end{document}
